% !TEX root = ../oddp.tex

\section{Interior simplices and the Milnor resolution}\label{s:milnor}
\begin{warning}
Recall from Section \cref{s:simplices} and its three warnings that we are not following the usual convention to denote the simplex category.
\end{warning}
Let $X, Y$ be pointed simplicial sets and let $X_+, Y_+$ be the augmented simplicial sets obtained by adding the trivial augmentation to a point $\{\mathrm{pt}\}$.

Let $X_+*_{\asimplex}Y_+$ denote the augmented simplicial join of $X_+$ with $Y_+$ and $X *_{\simplex} Y$ the non-augmented simplicial join of $X$ with $Y$. Let $|X|*|Y|$ denote the join of the realizations of $X$ and $Y$ and let $|X|\times |Y|$ denote their product. 

There are homeomorphisms\footnote{The second one can be explained as follows: $|X*_\simplex X|$ is obtained from $|X_+*_{\asimplex} X_+|$ as the result of collapsing the simplices in the image of $X*\{\mathrm{pt}\}$ and $\{\mathrm{pt}\}*X$ to the basepoint. These are precisely the canonical inclusions of $X$ into the join. The resulting space can be also obtained by adding two cones to the subspace $\{[(x,y,t)]\in X*X\mid t=1/2\}\cong X\times Y$}
\begin{align*}
	|X_+*_{\asimplex} Y_+| 
	& \cong 
	(|X| * |Y|)_+
	\simeq
	\Sigma(|X|\wedge |Y|)
	\\
	|X*_{\simplex} Y| & \cong \Sigma (|X|\times |Y|)
\end{align*}
and isomorphisms (recall our degree convention for the chain complex of an (augmented) simplicial set)
\begin{align} \label{eq:chains_join}
	\chains(X_+*_{\asimplex} Y_+) &\cong \chains(X_+)\ot \chains(Y_+)
	\\
	\chains(X*_{\simplex} Y) &\cong \chains(X)\ot \chains(Y).
\end{align}
and epimorphisms
\begin{align*}
	|X_+*_{\asimplex} Y_+| &\lra |X*_{\simplex} Y|
	\\
	\chains(X_+)\ot \chains(Y_+) & \lra
	\chains(X)\ot \chains(Y)	
\end{align*}

If $G$ is a finite group and $X$ is a simplicial set with a free $G$-action, then the infinite (augmented or non-augmented) join of $X$ with itself is a contractible space with a free $G$-action, and therefore a model of $EG$. Its normalised chain complex with coefficients in $R$ is then an acyclic resolution of $G$.

For example, taking $X$ to be the set $G$ with the left action by $G$ and using the non-augmented join we obtain the Milgram resolution of $G$. Taking instead the augmented join we obtain the Milnor resolution of $G$ (see \cref{ss:milgram}). Thus, we will refer to these kind of resolutions as \emph{Milgram} and \emph{Milnor} resolutions respectively.

If instead we take the $n$-fold (non-augmented) join of $X$ with itself and $X$ is $m$-connected, the resulting (non-augmented) simplicial set is $(n(m+2)-1)$-connected (resp. $(n+m)$-connected), hence its chain complex gives a partial resolution of $G$. We will be interested in the case $X=\partial \simplex^{r-1}$ is the boundary of the standard $(r-1)$-simplex, with the $\Cyc_r$-action that permutes the vertices of $\partial \simplex^{r-1}$. This action is free if and only if $r$ is prime. Write $\sigma = [0,1,\ldots,r-1]$ for the top simplex of $\asimplex^{r-1}$. The following definition assumes that $r$ is odd or that $R=\bF_2$.

\begin{definition}\label{def:milnor}
	The sequence of $R\Cyc_r$-chain complexes $[n]\mapsto \sus{n(1-r)}\chains(\partial\asimplex^{r})^{\ot n}$ has a cosemi-simplicial structure with coface maps
	\begin{align*}
		\d^i\colon \chains(\partial \asimplex^{r})^{\ot n} &\lra \sus{1-r}\chains(\partial \asimplex^{r})^{\ot n+1}
	\end{align*}
	given by $\d^i([\tau_0\ot \dots \ot\tau_{n-1}]) = \tau_0\ot \dots \tau_{i-1}\ot \partial\sigma \ot \rho^{-1}(\tau_i) \ot \dots \ot \rho^{-1}(\tau_{n-1}).
	$%\label{eq:simp_str2}
\end{definition}

Relabeling the vertices gives a simplicial homeomorphism
\[
\simplex^n*\overset{r}{\cdots}*\simplex^n
\lra
\simplex^{rn}
\lra
\simplex^{r}*\overset{n}{\cdots}*\simplex^{r}
\]
and using the isomorphisms in \eqref{eq:chains_join}, we obtain an isomorphism of chain complexes
\[
\alpha\colon \chains(\asimplex^n)^{\otimes r}
\to
\chains(\asimplex^{r})^{\otimes n}.
\]
Explicitly, if $\tau_0\ot \ldots \ot \tau_{r-1}$ is a generator of $\chains(\asimplex^{n})^{\ot r}$, and $\omega_i = [j \mid i\in \tau_j]$, then
\[
\alpha(\tau_0\ot \ldots \ot \tau_{r-1})
= (-1)^{s}\omega_0\ot\ldots \ot \omega_{r-1}.
\]
where $s$ is the sign of the permutation that rearranges the vertices.

Now, let $\beta$ be the automorphism of $\chains(\partial\asimplex^{r-1})^{\ot n}$ given by $\rho^0\ot \rho^{-1}\ot\ldots \ot \rho^{-n-1}$, where $\rho$ denotes the action of the cyclic group $\Cyc_r$. 

\begin{proposition}\label{lemma:bartomilnor}
	The restriction of $\beta\circ \alpha$ to the subcomplexes  $\chains(\asimplex^{\bullet})^{\ot r}_\interior$ induces an isomorphism of cosemi-simplicial chain complexes
	\[
	\sus{\bullet(1-r)}\chains(\asimplex^{\bullet})^{\ot r}_\interior
	\lra 
	\sus{\bullet(1-r)}\chains(\partial\asimplex^{r})^{\ot \bullet}.
	\]
\end{proposition}
\begin{proof}
	TBD: I think it is easier to write down the cosemi-simplicial structure in $\chains(\asimplex^{n})^{\ot r}_\interior$, see that $\alpha$ induces an isomorphism and that $\alpha$ commutes with the simplicial structure. \federico{TBD}
\end{proof}
